\section{An independent positive-block proof of the fold}
\label{sec:fold}

Let $D$ be the exact Boundary239 predecessor and $N$ the new table,
both with inherited labels. Put
\[
 \phi(169)=80,\qquad \phi(q)=q\text{ for }q\ne169,
 \qquad
 \Phi(q,h,T)=(\phi(q),h,T).
\]
The map fixes the start, dispatcher, and halt outcome.
Only old configurations need satisfy the old-state invariant;
new intermediate helper configurations may have other contexts.

\subsection{A structural invariant on arbitrary tapes}

Write $t_i=T(h+i)$. Define the predicate $P$ by the following
conditional requirements:
\[
 \begin{array}{c|l}
 q & \text{Required bits}\\ \hline
 67 & t_{-1}=0\\
 80 & (t_{-2},t_{-1})=(0,1)\\
 91 & (t_{-3},t_{-2},t_{-1})=(0,1,1)\\
 169 & t_{-1}=0.
 \end{array}
\]
Unlisted states impose no condition.

\begin{lemma}[Incoming-row invariant]
$P$ holds throughout every execution of $D$ starting at state 95,
for an arbitrary initial binary tape.
\end{lemma}
\begin{proof}
It is vacuous at state 95. The complete incoming inventory is:
states 52 and 124, reading zero, write zero and move right into 67;
state 67, reading one, writes one and moves right into 80;
state 80, reading one, writes one and moves right into 91; and
state 147, reading zero, writes zero and moves right into 169.
Each required nearest-left bit is the incoming write, and the
predecessor's requirements supply the older left bits. All other
targets have no requirement. The independent checker inspects
every predecessor row, so this inventory is verified for $D$ itself.
Induction on the execution length proves the result.
\end{proof}

One extra guard is needed:
\begin{equation}
 Q(q,h,T):\quad q=169\ \land\ T(h)=1
       \ \Longrightarrow\ T(h+1)=1.
 \label{eq:nextone}
\end{equation}
Unlike $P$, we claim $Q$ only on the actual blank initialization and
the valid source-generated runtime, rather than on every initial
tape.

\subsection{Four exact local paths}

The four exceptional old rows and their new state paths are
\begin{align*}
 (80,0):\quad&
 80\to106\to92\to107\to125\to25
 &&(5\text{ steps}),\\
 (169,0):\quad&
 80\to106\to93\to57
 &&(3\text{ steps}),\\
 (169,1):\quad&
 80\to91\to167\to151\to173\to56
 &&(5\text{ steps}),\\
 (91,1):\quad&
 91\to167\to151\to24\to36\to54\\
 &\hspace{7mm}\to71\to82\to94\to50
 &&(9\text{ steps}).
\end{align*}
The final two lines describe one nine-transition block.

\begin{lemma}[Four whole-tape identities]
\label{lem:fourblocks}
Under $P$ and, for the third path, $Q$, each displayed new path has
the same final tape and head as the corresponding single old step,
and its final state is the image of the old final state under $\phi$.
Every other one of the 474 old rows commutes with $\Phi$ in one step.
\end{lemma}
\begin{proof}
For the first path the required footprint is offsets $-1,0$ with
bits $1,0$; for the second it is $-1,0$ with bits $0,0$.
The third requires $(-1,0,1)=(0,1,1)$, while the fourth requires
$(-3,-2,-1,0)=(0,1,1,1)$. Direct substitution into the literal table
gives the displayed paths. The first two preserve the tape and
end one cell right. The third erases the original current one
and ends one cell right. The fourth erases the original current
one and ends one cell left.

The independent verifier assigns a distinct symbolic original
bit to every unspecified cell. An attempted read of an
unconstrained bit is rejected. It compares the two complete
finite write maps with the same symbolic exterior, along with
head position and mapped state. Thus the four comparisons hold
for arbitrary tapes outside their footprints. There is no internal
halt. For all other predecessor rows the mapped literal row is
checked directly, including the row entering halt 254.
\end{proof}

\subsection{The extra guard survives the source change}

The observer \file{verification/verify_guard238.py} instruments
the unchanged affine/list engine while it proves \emph{Boundary239's}
runtime. At every state-169 entry, it checks a rightward scan into
an explicit run. When that run is a one-run, the constant length
must be at least two.

All 8,386 interfaces pass, using 2,247,513 symbolic rules and all
478 predecessor rows. The 32 explicit state-169/read-one visits
have run lengths
\[
 g+K,\qquad g\ge0,\qquad 23\le K\le59.
\]
This proves the neighbor in \eqref{eq:nextone} is one. The value
23 is the new minimum for the changed source; the earlier
Rotation240 guard's value 21 is not used.

No accelerated rule hides an unchecked visit. State 169 is
neither an entry of a list schema nor a source or target inside
any of the four explicit one-block proofs. Neither of its rows
is a self-loop, so a unary acceleration cannot conceal such a
visit either. The exact predecessor blank prefix has six
state-169/read-one visits, each with the required right neighbor.
Together with the source/frame composition, these facts prove
$Q$ throughout the predecessor's actual blank run.

\begin{theorem}[Guarded composition]
\label{thm:fold}
$D$ and $N$ have equivalent halting behavior on the blank tape.
If $D$ performs $n$ transitions before a checkpoint, then $N$
reaches the mapped checkpoint after
\begin{align}
 \tau(n)=n
 &+4V_{80,0}(n)+2V_{169,0}(n)\notag\\
 &+4V_{169,1}(n)+8V_{91,1}(n), \label{eq:clock}
\end{align}
where $V_{q,b}(n)$ counts the indicated old row among the first
$n$ transitions. In particular, $n\le\tau(n)\le9n$.
\end{theorem}
\begin{proof}
The invariant and guard just proved permit
Lemma~\ref{lem:fourblocks} at every old configuration. Inductively
concatenate the applicable new blocks. Equation~\eqref{eq:clock}
counts their extra transitions, and every block has length
between one and nine.

If the old machine halts, the corresponding new checkpoint is
halted. If it runs forever, the strictly increasing checkpoint
times are unbounded; no block halts internally. Determinism
identifies this constructed trajectory with the new machine's
unique trajectory. Hence the new machine also runs forever.
This establishes both directions of halting equivalence.
\end{proof}

\subsection{Independent initialization clock}

The local verifier compares all 409,692 full-tape checkpoints
spanning the predecessor's 409,691-transition initialization.
The exceptional old-row visits, in the order used in
\eqref{eq:clock}, are
\[
 (18,12,6,3).
\]
Therefore
\[
 \tau(409691)
 =409691+4\cdot18+2\cdot12+4\cdot6+8\cdot3
 =409835.
\]
Both endpoints are state 96, head $-5$, and the same 351 one-cells.
This independently agrees with the direct sparse, dense, and
canonically relabeled startup check.

The useful general principle is that two verified transformations
compose once the second transformation's configuration obligations
are proved for the first transformation's actual output system.
Disjoint text edits are only a discovery heuristic. Here the
positive local simulation and the re-established guard supply the
missing mathematical compatibility condition.
